Table~\ref{tab:rewrite-rsi-accuracy} reports the final numerical checks.
\begin{table}[H]
 \centering
 \caption{Numerical accuracy of the RSI-activation equilibria}
 \label{tab:rewrite-rsi-accuracy}
 \small
 \begin{tabular}{rrrrr}
 \toprule
 $\sigma$ & Final horizon & Mesh points & Dynamic residual & Horizon change\\
 \midrule
 0.90 & 4,815.9 & 1,740 & $1.63\times10^{-9}$ & $2.89\times10^{-8}$\\
 1.00 & 4,591.9 & 1,868 & $2.84\times10^{-9}$ & $3.76\times10^{-8}$\\
 1.10 & 4,591.9 & 1,959 & $3.18\times10^{-9}$ & $2.86\times10^{-8}$\\
 1.50 & 2,545.7 & 2,294 & $3.65\times10^{-9}$ & $1.52\times10^{-8}$\\
 \bottomrule
 \end{tabular}
 \par\smallskip
 \begin{minipage}{0.95\textwidth}\footnotesize
 Dynamic residuals are the largest independently reconstructed errors
 across the full-horizon and dense first-decade checks. Horizon change
 compares the last two solutions in detrended logarithmic coordinates
 over years 0--500. Error bounds are rounded upward. The event-continuity
 log gaps are below $1.4\times10^{-15}$; the final unit-elastic price ratio
 is $0.2000000012$.
 \end{minipage}
\end{table}
